\documentclass[a4paper]{article}
\usepackage[pdftex]{graphicx}
\usepackage[utf8]{inputenc}
\usepackage{enumerate}
\usepackage{amssymb}
\usepackage{colortbl}
\usepackage{icomma}
\usepackage{colortbl}
\usepackage{musixtex}
\usepackage{siunitx}
\sisetup{locale=DE}
\usepackage{geometry}
\geometry{a4paper, top=15mm, left=15mm, right=15mm, bottom=15mm, 
	headsep=10mm, footskip=12mm}
\usepackage{href-ul}
\hypersetup{ 
	colorlinks=true, 
	linkcolor=blue, 
	urlcolor=blue}
\begin{document}
	
{\bf The scale, mathematically seen}
	
	\begin{enumerate}[1.]
	
{\bf \item	Fill in the table using the formula $y = 440 \cdot 2^{\frac{x}{12}}$ }.

Here, x is the number of semitones and y is the frequency.

\renewcommand{\arraystretch}{2}
\begin{tabular}{|>{\columncolor[gray]{.8}}p{1.5 cm}|p{1.5 cm}|p{1.5 cm}|p{1.5 cm}|p{1.5 cm}|p{1.5 cm}|p{1.5 cm}|p{1.5 cm}|} 
	\hline 
	\rowcolor[gray]{.8} tone & a, & c & c\# & d & d\# & e & f\\ 
	\hline 
	x & -12 & -9 & -8 &-7 & -6& -5 & -4 \\ 
	\hline 
	y in Hz & & & & & & & \\ 
	\hline
\end{tabular}


\renewcommand{\arraystretch}{2}
\begin{tabular}{|>{\columncolor[gray]{.8}}p{1.5 cm}|p{1.5 cm}|p{1.5 cm}|p{1.5 cm}|p{1.5 cm}|p{1.5 cm}|p{1.5 cm}|p{1.5 cm}|p{1.5 cm}|} 
	\hline 
	\rowcolor[gray]{.8} tone & f\# & g & g\# & a & a\# & h & c' & a'\\ 
	\hline 
	x & -3 & -2 & -1 & 0 & 1 & 2 & 3 & 12 \\ 
	\hline 
	y in Hz & & & & 440 & & & &\\ 
	\hline
\end{tabular}


{\bf \item Graph the frequencies over the semitone steps}:\\
{\bf a) with linear y-axis}\\
\begin{center}
	\vspace{-0,5 cm}
	\includegraphics[width=15 cm]{expto01.png} \\
	\vspace{-0,5 cm}
	\hspace{0,0 cm}
	\begin{music} 
		\parindent10mm 
		\instrumentnumber{1} 
		\setstaffs1{1} 
		\startextract 
		\Notes 
		\qa {a}
		\hskip 0,5  cm
		\qa c
		\qa{^c}	
		\qa d  
		\qa{^d}
		\qa e 
		\qa f 
		\qa{^f}
		\qa {g}
		\qa{^g'}	
		\qa{a}
		\qa{^a}	
		\qa b 
		\qa {c'}
		\hskip 2 cm
		\qa{a''} 
		\en 
		\endextract
	\end{music}
\end{center}
{\bf	b) with logarithmic y-axis}\\
	
\begin{center}
	\vspace{-0.5 cm}
	\includegraphics[width=15 cm]{logto01.png} \\
\end{center}

\begin{music} 
	\parindent10mm 
	\instrumentnumber{1} 
	\setstaffs1{1} 
	\startextract 
	\Notes 
	\qa c 
	\qa{^c} 
	\qa d 
	\qa{^d} 
	\qa e 
	\qa f 
	\qa{^f} 
	\qa g 
	\qa{^g} 
	\qa {'a}
	\qa{^a}
	\qa b
	\qa {c'}
	
	\en
	\endextract
\end{music}
{\bf \item Fill in the table with the musical intervals}

\renewcommand{\arraystretch}{1}
\begin{tabular}{|>{\columncolor[gray]{.8}}p{3 cm}|p{2.5 cm}|p{3 cm}|p{2.5 cm}|p{4 cm}|}
	\hline
	\rowcolor[gray]{.8} Name & Semitone Steps & Notes & $2^{x/12}$ & $\approx$ Frequency Ratio *\\
	
	\hline
	Prime & &
	\begin{music}
		\instrumentnumber{1}
		\setstaffs1{1}
		
		\startextract
		\Notes \zcharnote{c}{\raisebox{-3ex}{C}} \qa c \hskip 1.2cm \en
		
		\endextract
	\end{music}
	
	& &\\
	
	\hline
	minor second & 1 &
\begin{music}
	\instrumentnumber{1}
	\setstaffs1{1}
	\startextract
	\Notes \zcharnote{c}{\raisebox{-3ex}{C}} \qa c \qa{^c} \en
	\endextract
\end{music}

& 1,059& 16:15 \\
	
	\hline
	major second & &
	\begin{music}
		\instrumentnumber{1}
		\setstaffs1{1}
		
		\startextract
		\Notes \zcharnote{c}{\raisebox{-3ex}{C}} \qa c \hskip 1.2cm \en 
		\endextract 
	\end{music} 
	& & \\ 
	\hline 
	minor third & & 
	\begin{music} 
		\instrumentnumber{1} 
		\setstaffs1{1} 
		\startextract 
		\Notes \zcharnote{c}{\raisebox{-3ex}{C}} \qa c \hskip 1.2cm \en 
		\endextract 
	\end{music} 
	& & \\ 
	\hline 
	major third & & 
	\begin{music} 
		\instrumentnumber{1} 
		\setstaffs1{1} 
		\startextract 
		\Notes \zcharnote{c}{\raisebox{-3ex}{C}} \qa c \hskip 1.2cm \en 
		\endextract
	\end{music}
	
	& & \\
	
	\hline
	Fourth & &
	
	\begin{music}
		\instrumentnumber{1}
		\setstaffs1{1}
		
		\startextract
		\Notes \zcharnote{c}{\raisebox{-3ex}{C}} \qa c \hskip 1.2cm \en
		
		\endextract
	\end{music}
	
	& & \\
	
	\hline
\end{tabular}

\vspace{0.5 cm}

	* There are differences between uniform and natural tuning. The approximate values for the frequency ratios are:

\quad \fbox{1 : 1}
\quad \fbox{5:4}
\quad\fbox{4:3}
\quad\fbox{6:5}
\quad\fbox{9:8}

\newpage
\begin{music}
	\parindent10mm
	
	\instrumentnumber{1} 
	
	\setstaffs1{1} 
	\startextract 
	\Notes 
	\qa c 
	\qa{^c} 
	\qa d 
	\qa{^d} 
	\qa e 
	\qa f 
	\qa{^f} 
	\qa g 
	\qa{^g} 
	\qa {'a} 
	\qa{^a} 
	\qa b 
	\qa {c'} 
	\en 
	\endextract
\end{music}
\renewcommand{\arraystretch}{1}
\begin{tabular}{|>{\columncolor[gray]{.8}}p{3 cm}|p{2.5 cm}|p{3 cm}|p{2.5 cm}|p{3 cm}|} 
	\hline 
	\rowcolor[gray]{.8} Name & Semitone Steps & Notes & $2^{x/12}$ & $\approx$ Frequency Ratio *\\ 
	\hline 
	tritone & & 
	\begin{music} 
		\instrumentnumber{1} 
		\setstaffs1{1} 
		\startextract 
		\Notes \zcharnote{c}{\raisebox{-3ex}{C}} \qa c \hskip 1.2cm \en 
		\endextract 
	\end{music} 
	& & \\ 
	\hline 
	fifth & & 
	\begin{music} 
		\instrumentnumber{1} 
		\setstaffs1{1} 
		\startextract 
		\Notes \zcharnote{c}{\raisebox{-3ex}{C}} \qa c \hskip 1.2cm \en 
		\endextract 
	\end{music} 
	& & \\ 
	\hline 
	minor sixth & & 
	\begin{music} 
		\instrumentnumber{1} 
		\setstaffs1{1} 
		\startextract 
		\Notes \zcharnote{c}{\raisebox{-3ex}{C}} \qa c \hskip 1.2cm \en 
		\endextract 
	\end{music}
	& & \\
	
	\hline
	major sixth & &
	
	\begin{music}
		\instrumentnumber{1}
		\setstaffs1{1}
		
		\startextract
		
		\Notes \zcharnote{c}{\raisebox{-3ex}{C}} \qa c \hskip 1.2cm \en
		
		\endextract
	\end{music}
	
	& & \\
	
	\hline
	minor seventh & &
	
	\begin{music}
		
		\instrumentnumber{1}
		\setstaffs1{1}
		
		\startextract
		
		\Notes \zcharnote{c}{\raisebox{-3ex}{C}} \qa c \hskip 1.2cm \en
		
		\endextract
	\end{music}
	
	& & \\
	
	\hline
	major seventh & &
	
	\begin{music} 
		
		\instrumentnumber{1} 
		\setstaffs1{1} 
		\startextract 
		\Notes \zcharnote{c}{\raisebox{-3ex}{C}} \qa c \hskip 1.2cm \en 
		\endextract 
	\end{music} 
	& & \\ 
	\hline 
	octave & & 
	\begin{music} 
		\instrumentnumber{1} 
		\setstaffs1{1} 
		\startextract 
		\Notes \zcharnote{c}{\raisebox{-3ex}{C}} \qa c \hskip 1.2cm \en 
		\endextract 
	\end{music} 
	& & \\ 
	\hline
	\end{tabular}
	
	\vspace{0.5 cm}
	* There are differences between uniform and natural tuning. The approximate values for the frequency ratios are:

	\quad \fbox{3:2}
	\quad \fbox{5:3}
	\quad \fbox{45:32}
	\quad \fbox{2:1}
	\quad \fbox{9:5}
	\quad \fbox{8:5}
	\quad \fbox{15:8}
	
	\vspace{1 cm}
	{\bf \item What does a semitone step have to do with $\sqrt[12]{2}$?}
	\vspace{1 cm}
	{\bf \item Play the intervals with a piano-app. Which one is your favorite?}.
	
	\end{enumerate} 
	
	\fbox{
	\begin{minipage}{0.5\textwidth}
		For the solution, please \href{https://www.okuyakl.de/math/m10tonL135/le135.pdf}{click here} or scan the QR-Code.\\
		More worksheets are available at
		
		\href{https://www.okuyakl.de}{www.okuyakl.de}
	\end{minipage}
	\hfill
	\begin{minipage}{0.4\textwidth}
		\includegraphics[width=1.5 cm]{../../viecher/zwe03}
		\includegraphics[width=3 cm]{qre135}
		\includegraphics[width=2 cm]{../../viecher/afanticon1}
		
		\end{minipage}}
		
\end{document}
